\section{Weighted configurations and critical exponents}
\label{sec:critical}

The geometric problem is to bound the average multiplicity of shaded tubes
in terms of their shading density and concentration in planks. We first
express this problem for finite weighted, indexed families of affine lines.
Restrictions and changes of scale preserve inherited weights, but may
produce repeated line parameters, so the class includes those repetitions.
We then define a family of critical inequalities and construct sequences
that attain their optimal exponents. These sequences will be the input to
the local geometric arguments.

\subsection{The indexed model and logarithmic comparisons}

Let $N=2^j$ tend to infinity.  An index $i$ consists of a positive weight
$\omega_i$, a matrix $M_i=(b_i,u_i)$ in a fixed bounded subset of
$\R^{2\times2}$, and a nonempty subset $S_i$ of the $N$ dyadic time bins in
$[0,1]$.  Different indices may have identical matrices and different
shadings.  The spatial trace is
\[
 M_i(t)=b_i+tu_i.
\]
An event is a pair $(i,J)$ with $J\in S_i$.  Its location is
$(t_J,M_i(t_J))$, where $t_J$ is the bin center.  Let $E$ be the union of the
space--time grid cells of side $N^{-1}$ containing these locations.  Write
\begin{equation}\label{eq:indexed-incidence}
 n=\sum_i\omega_i,
 \qquad \mathcal I=\sum_i\omega_i|S_i|,
 \qquad k=\mathcal I/n,
 \qquad m=\mathcal I/|E|.
\end{equation}
Thus $k$ is initially a weighted mean, and $m$ is the raw average
multiplicity.  After regularization all $|S_i|$ will have the same exponent
as $k$.  Each event has mass $\omega_i$; probabilities on the event set
always mean division by $\mathcal I$.

A full-time plank of widths $a,b$, where $1\le a\le b\le N$, is a set with
cross-sections
\[
 \left\{x\in\R^2:
 |(x-b_0-tu_0)\cdot e_1|\le a/N,
 \quad |(x-b_0-tu_0)\cdot e_2|\le b/N\right\},
 \qquad 0\le t\le1,
\]
for an orthonormal spatial frame $(e_1,e_2)$ and affine center $b_0+tu_0$.
Constant factors in these widths are permitted.  An index belongs to a
plank only if its entire trace over the indicated time interval is
contained in it.  We use the same definition on shorter dyadic time blocks;
there the admissible widths are at most the block length in finest-cell
units.  In particular, membership is stronger than having an event in the
plank.  The mass $n_B$ of a plank is the sum of the inherited weights of its
contained indices.

Fixed enlargements of cells, moving event locations a bounded number of
fine cells, and replacing the bounded matrix patch by any other fixed-size
patch do not affect the exponents below.  For the last assertion, translate
the matrix center and rescale space by a fixed constant.  Both the old and
new support grids admit bounded-overlap coverings by the other grid, and
plank tests transfer after bounded enlargement and a fixed finite covering.

\begin{convention}[Exponent notation]\label{conv:exponents}
For positive quantities on a sequence with $N\to\infty$, write
\[
 A\leexp B
 \quad\Longleftrightarrow\quad
 \limsup_{N\to\infty}\frac{\log(A/B)}{\log N}\le0,
 \qquad
 A\eqexp B
 \quad\Longleftrightarrow\quad
 \frac{\log(A/B)}{\log N}\longrightarrow0.
\]
A factor $N^{o(1)}$ has logarithm $o(\log N)$, uniformly over the objects
being compared in that statement.  Subsequence selection is allowed.
Every contradiction will compare exponents separated by a fixed positive
amount before the resolution tends to infinity.  All these comparisons are
homogeneous in the weights.
\end{convention}

A refinement deletes events and then deletes empty indices; retained
indices keep their original weights.  It is \emph{extensive} if its event
mass is at least $N^{-o(1)}\mathcal I$.  A restriction of loss exponent at
most $e$ retains at least $N^{-e+o(1)}\mathcal I$.

\begin{definition}[Uniform temporal profile]\label{def:uniform-profile}
An admissible profile is a nondecreasing $1$-Lipschitz function
$F:[0,1]\to[0,1]$ with $F(0)=0$.  A sequence of indexed configurations has
uniform profile $F$ if, uniformly over every retained index and every
occupied dyadic time block containing $K$ fine bins,
\begin{equation}\label{eq:uniform-profile}
 |S_i\cap I|=N^{\kappa-F(\kappa)+o(1)},
 \qquad \kappa=\log_N K.
\end{equation}
The uniform error also applies over the scale range $0\le\kappa\le1$.
In this case $k\eqexp N^{1-\alpha}$, where $\alpha=F(1)$.
\end{definition}

The variable $\kappa$ measures block size, not physical time: increasing
$\kappa$ moves from fine bins to coarser time blocks. A block at this scale
has $N^\kappa$ available bins, so $F(\kappa)$ records the exponent lost
from full occupation. Between scales $s<t$, the number of occupied
$s$-blocks in an occupied $t$-block has exponent
$(t-s)-(F(t)-F(s))$.

Complete shading has $F=0$, while one marked bin per index gives
$F(\kappa)=\kappa$. Total density alone does not determine the profile.
For example, for $0<\alpha<1$ the admissible profiles
\[
 F_1(\kappa)=\alpha\kappa,
 \qquad F_2(\kappa)=(\kappa-1+\alpha)_+
\]
have the same endpoint deficit $\alpha$; here $r_+=\max\{r,0\}$.
The first distributes the deficit
at a constant rate across scales. The second is realized by marking one
aligned dyadic interval of $N^{1-\alpha+o(1)}$ bins on every index:
occupied blocks below that interval's size are full, and larger blocks
contain the entire marked interval. This distinction matters when we
change scales or fill occupied blocks.

It is enough to arrange~\eqref{eq:uniform-profile} on meshes whose maximal
gap tends to zero sufficiently slowly.  Indeed the count in an occupied
intermediate block is bounded below by that in one of its occupied finer
mesh blocks and above by that in its coarser mesh ancestor.  The two mesh
exponents differ by at most their scale gap plus the errors at the mesh
points.

\subsection{The density bound and its deformation}

Put $\lambda=k/N$. The multiplicity estimate we seek, for $p>2$
arbitrarily close to two, is
\begin{equation}\label{eq:undeformed-density-target}
 m\leexp \lambda^{-p}\sup_B\frac{n_B}{ab}.
\end{equation}
Here $B$ ranges over full-time planks of widths $a,b$. To see the
connection with the maximal theorem, consider unit-weight lines with
pairwise $c/N$-separated slopes, where $c>0$ is fixed. Comparing the two
endpoint cross-sections of a plank confines those slopes to a rotated
rectangle of side lengths $O(a/N)$ and $O(b/N)$. Slope separation therefore
gives $n_B\le C ab$. Since $m=nk/|E|$, the desired bound would give
\[
 |E|\geexp nN\lambda^{p+1}.
\]
As $p$ decreases to two, this approaches the cubic density bound required
for the $L^3$ maximal estimate. The final passage from this discrete
estimate to arbitrary functions is proved in
Proposition~\ref{prop:maximal-reduction}.

We deform both the temporal penalty and the powers assigned to the two
plank widths. The resulting critical exponent measures the additional
power of resolution needed in the bound. One endpoint of the deformation
will be elementary; derivatives of the optimal exponent at an interior
point will later select temporal deficits and plank shapes.

Fix $0<\eta<1/2$.  For $p>2$, $z=-q\in[-1,0]$ and $d=2-q$, set
\begin{equation}\label{eq:temporal-integrand}
 \mathcal D(q,e)=\frac1\eta\int_{1/2-\eta}^{1/2}\min(q,re)\,dr,
 \qquad P_{p,q}(e)=pe-\mathcal D(q,e),\qquad 0\le e\le1.
\end{equation}
For an admissible profile define
\begin{equation}\label{eq:temporal-cost}
 \Pi_F(t)=\int_0^t P_{p,q}(F'(u))\,du,
 \qquad 0\le t\le1.
\end{equation}
The derivative exists almost everywhere because $F$ is Lipschitz.  When the
profile and parameters are fixed we write $\Pi$ for $\Pi_F$.  Dependence on
the fixed smoothing width $\eta$ is suppressed. At $q=0$,
$\Pi_F(1)=pF(1)$, so only the total density enters. For $q>0$ the
discount $\mathcal D(q,F')$ also records how the deficit is distributed
across scales. The averaging in $r$ smooths this discount in $q$, as needed for
the later derivative test.

For fixed $q$, the function $e\mapsto\min(q,re)$ is increasing, concave,
and $r$-Lipschitz. Thus $P_{p,q}$ is increasing and convex, vanishes at
zero, and has all secant slopes in $[p-1/2,p]$. The bounds
$0\le\mathcal D(q,e)\le\min(q,e/2)$ also give
\begin{equation}\label{eq:cost-elementary-bounds}
 (p-1/2)F(t)\le\Pi_F(t)\le pF(t),
 \qquad \Pi_F(1)\ge pF(1)-q.
\end{equation}
These elementary estimates suffice to define the critical problem and
show that its admissible exponents form a nonempty set.

For $0\le A\le d-1$, define the plank parameter
\begin{equation}\label{eq:plank-parameter}
 \Delta_A=\sup_B\frac{n_B}{a^{d-1-A}b^{1+A}},
\end{equation}
where $B$ runs over full-time planks.  The letter $A$ in this definition is
only a candidate exponent.  Every configuration is finite, and the
supremum may equivalently be taken over a finite net with bounded
thickening.  The bounds below are independent of the size of that net.

\begin{definition}[Critical exponent]\label{def:critical}
For fixed $p,z,\eta$, let $h(p,z)$ be the least $A\in[0,d-1]$ such that
for every sequence with a uniform profile $F$,
\begin{equation}\label{eq:critical-bound}
 m\leexp N^{A+\Pi_F(1)}\Delta_A.
\end{equation}
\end{definition}

The definition ranges over all bounded-patch weighted indexed
configurations. At $q=0$ one has $d=2$, and the candidate $A=0$ gives
exactly \eqref{eq:undeformed-density-target} for a uniform profile.
There is then no dependence on $\eta$. Our task is to show that
$h(p,0)=0$ for $p>2$ arbitrarily close to two.

More precisely, Proposition~\ref{prop:maximal-reduction} shows that
$h(p,0)=0$ implies, for every $\varepsilon>0$,
\[
 \|K_\delta f\|_{L^3(S^2)}
 \le C_{p,\varepsilon}\delta^{-(p-2)/3-\varepsilon}
          \|f\|_{L^3(\R^3)}.
\]

\begin{lemma}[Existence and elementary comparisons]
\label{lem:critical-existence}
Definition~\ref{def:critical} is well-defined.  Its admissible candidates
form a closed upper interval in $[0,d-1]$, and $h(p,-1)=0$.  For
every $0\le A_1\le A_2\le d-1$,
\begin{equation}\label{eq:plank-exponent-Lipschitz}
 N^{-(A_2-A_1)}\Delta_{A_1}
 \le\Delta_{A_2}\le\Delta_{A_1},
 \qquad
 N^{A_1}\Delta_{A_1}\le N^{A_2}\Delta_{A_2}
 \le N^{A_2-A_1}N^{A_1}\Delta_{A_1}.
\end{equation}
All fixed geometric constants may be absorbed in exponent notation.
\end{lemma}

\begin{proof}
The comparisons follow by multiplying each plank denominator by
$(b/a)^{A_2-A_1}$ and using $1\le b/a\le N$.  Thus a fixed plank's
contribution to $N^A\Delta_A$ has logarithmic slope
$\log_N(Na/b)\in[0,1]$.  Taking the supremum preserves these comparisons.
Monotonicity of the bound and closedness of the set of valid candidates
follow, the latter by letting $A_2-A_1$ tend to zero after applying the bound
on an arbitrary fixed sequence.

It remains to prove validity at $A=d-1$. We bound multiplicity by the
weighted count of pairs meeting in fine cells, then sum by slope separation.
Let $\mu(Q)$ be the total weight
of events in a fine space--time cell $Q$.  Cauchy--Schwarz gives
\[
 m=\mathcal I/|E|\le
 \frac{\sum_Q\mu(Q)^2}{\mathcal I}.
\]
Fix one indexed line.  Other lines with relative slope magnitude comparable
to $\theta\ge N^{-1}$ that pass within bounded fine distance of it at some
time are contained in a full-time plank of both widths $O(\theta N)$
around it.  Their total mass is at most
$O(\Delta_{d-1}(\theta N)^d)$.  Two such lines can be within bounded fine
distance for at most $O(1/\theta)$ time bins: their spatial difference is
affine with velocity of magnitude comparable to $\theta$.  For relative
slopes at most $N^{-1}$, use mass $O(\Delta_{d-1})$ and at most $N$ bins.
Fixed large slope ranges or patch boundaries require only fixed finite
coverings.  Summing over dyadic $\theta$ and then the reference line gives
\[
 \sum_Q\mu(Q)^2
 \lesssim n\Delta_{d-1}N^d
       \left(1+\sum_{\substack{\theta\text{ dyadic}\\N^{-1}\le\theta\lesssim1}}
                              \theta^{d-1}\right)
 \le n\Delta_{d-1}N^{d+o(1)}.
\]
This includes $d=1$, where the sum costs a logarithm, and includes repeated
indices.  Consequently
\[
 m\leexp\Delta_{d-1}N^d/k
 \eqexp\Delta_{d-1}N^{d-1+F(1)}.
\]
Since $\Pi_F(1)\ge(p-1/2)F(1)\ge F(1)$,
\eqref{eq:critical-bound} holds at $A=d-1$.  When $z=-1$, one has $d=1$
and this is the only possible candidate, proving $h(p,-1)=0$.
\end{proof}

\subsection{Finite cleaning and profile regularization}

To apply the critical inequality after selecting events, we need uniform
profiles and lower bounds on the mass retained in every tested cell.
The following elementary deletion estimate preserves those lower bounds.
Its partitions are allowed to have arbitrarily many cells; a bound on the number of cells is needed only for the probability
pigeonholing that follows it.

\begin{lemma}[Finite cleaning]\label{lem:finite-cleaning}
Let $\Omega$ be a finite set with positive measure $\mu$, let
$A\subseteq\Omega$ satisfy $\mu(A)\ge\rho\mu(\Omega)$, and let
$\mathcal P_1,\ldots,\mathcal P_L$ be partitions of $\Omega$.
If $0<\vartheta_0<\rho/L$, there is $A'\subseteq A$ with
\[
 \mu(A')\ge(\rho-L\vartheta_0)\mu(\Omega)
\]
such that every cell $C$ of a tested partition meeting $A'$ satisfies
\[
 \mu(A'\cap C)\ge\vartheta_0\mu(C).
\]
In particular one can take $\vartheta_0=\rho/(2L)$ and retain at least
$\rho\mu(\Omega)/2$.
\end{lemma}

\begin{proof}
Start with $A$.  Whenever a tested cell has positive remaining mass less
than $\vartheta_0\mu(C)$, delete its remaining points.  A cell used in such
a deletion becomes empty and is never charged again.  Each deletion costs
less than $\vartheta_0\mu(C)$, measured against the \emph{original} measure.
For each partition the sum of the original cell masses is $\mu(\Omega)$.
Thus all deletions together cost at most $L\vartheta_0\mu(\Omega)$.
The process terminates because $\Omega$ is finite.  Cascades between the
partitions cause no additional charge.
\end{proof}

A conditional probability test is specified by a partition into domains and
by a partition of each domain into at most $N^C$ cells, with $C$ fixed for
the test.  Its value at an event is the probability of that event's cell
conditional on its domain.  The number of domains need not be polynomial.
For example, the domains may be the individual indices.  For any $A>C$,
the events in conditional cells of probability below $N^{-A}$ have total
probability at most $N^{C-A}$, by summing within each domain and then
averaging over the domains.  A global partition with at most $N^C$ cells is
the special case of one domain.

For a fixed finite list of such tests, discard these negligible tails and
round each remaining $-\log_N$ conditional probability to an interval of
length $\nu$.  There are at most
\[
 \prod_{r=1}^L(2+A_r/\nu)
\]
possible label vectors.  One vector carries at least the reciprocal of this
number, up to the discarded mass.  Clean both the domain and cell partitions
using Lemma~\ref{lem:finite-cleaning}.  On a surviving cell, both numerator
and denominator of a conditional probability lie between their original
values and $\vartheta_0$ times those values.  Its conditional exponent
therefore changes by at most $|\log_N\vartheta_0|$.  This proves the stated
pigeonholing simultaneously for probabilities and conditional
probabilities.  Further finite partitions can be included solely for
cleaning, without pigeonholing their cell masses.

Within an index, temporal counts can be used instead of conditional
probabilities.  They are integers between $1$ and $N$ and their logarithmic
labels have a polynomially bounded range regardless of the number or
weights of indices.

\begin{lemma}[Regularization and inherited branching]
\label{lem:regularization}
The following assertions hold.
\begin{enumerate}[label=(\roman*)]
\item Every sequence of indexed configurations has, after passage to a
subsequence, an extensive refinement with a uniform temporal profile.  A
prescribed finite list of the preceding probability and conditional-count
tests can be regularized at the same time.  Countably many prescribed tests
can be imposed by the finite-test diagonal described below.
\item Suppose an original sequence has profile $F$, and a restriction
followed by extensive regularization has profile $G$.  Then
\begin{equation}\label{eq:branching-domination}
 G(t)-G(s)\ge F(t)-F(s),\qquad 0\le s\le t\le1.
\end{equation}
If the restriction retains at least $N^{-e+o(1)}$ of the original event
mass, then
\begin{equation}\label{eq:total-deficit-loss}
 0\le G(1)-F(1)\le e.
\end{equation}
In particular an extensive refinement preserves $F$.
\item The same branching comparison holds within a fixed interval of
logarithmic scales when events are routed through charts, provided each
derived temporal fiber is a subset of one original index and starting
block, its weight is inherited, and each original index and starting block
produces at most $N^{o(1)}$ chart fibers.  Total event mass and total
available starting-block weight are counted over that ensemble.  If the
routed ensemble retains at least $N^{-e+o(1)}$ of the original ensemble's
event mass and is extensively regularized to a common profile, its total
deficit increase on the interval is at most $e$ in original logarithmic
units.  Without that retained-mass hypothesis only the increment comparison
is asserted.
\end{enumerate}
\end{lemma}

\begin{proof}
Fix first a finite mesh
$0=\kappa_0<\cdots<\kappa_J=1$ and an accuracy $\nu>0$.
At resolution $N$, round each interior mesh point to a multiple of
$1/\log_2N$, keeping the endpoints. For sufficiently large $N$ these
rounded points remain ordered; they specify actual dyadic block sizes and
differ from the prescribed points by at most $1/\log_2N$. In this finite construction
$\kappa_r$ denotes the rounded value. The rounding error tends to zero
before the mesh is refined.
At every event label the original counts in its index's occupied ancestral
blocks at these scales, rounding $\log_N$ of each count to intervals of
length $\nu$.  There are at most $(2+1/\nu)^{J+1}$ labels.  Include the
other prescribed labels, retain a label class of maximal event mass, and
clean all the partitions into pairs consisting of an index and a block.
Those partitions may have many cells, which is permitted in
Lemma~\ref{lem:finite-cleaning}. Let $c_j$ be the lower endpoint, a
multiple of $\nu$, of the common rounded count interval at scale
$\kappa_j$. Every surviving occupied block there has count
between
\[
 \vartheta_0N^{c_j}
 \quad\hbox{and}\quad N^{c_j+\nu}.
\]
The same assertion at the root gives uniform counts per retained index.
The weights cancel in these within-index comparisons.

Put $\lambda_N=\log_N(1/\vartheta_0)$ and
$\epsilon_N=\nu+\lambda_N$.  An occupied coarse block contains an
occupied finer block.  Conversely it has at most $N^{\kappa_k-\kappa_j}$
subblocks across the actual dyadic scale gap $[\kappa_j,\kappa_k]$.
The preceding count bounds therefore give, for every $j<k$,
\[
 -\epsilon_N\le c_k-c_j
 \le\kappa_k-\kappa_j+\epsilon_N.
\]
The error does not accumulate with the number of intervening scales.
To replace these approximate count exponents by an admissible profile,
first define a branching-count exponent and then its deficit:
\[
 b_N(s)=\min_{0\le j\le J}\{c_j+(s-\kappa_j)_+\},
 \qquad F_N(s)=s-b_N(s).
\]
Each function in the minimum is nondecreasing and $1$-Lipschitz, and so is
their minimum.  Since $c_0=0$ and every $c_j\ge0$, one has
$0\le b_N(s)\le s$ and $b_N(0)=0$.  Thus $F_N$ is an admissible profile.
At a tested scale the preceding increment bounds imply
\[
 c_j-\epsilon_N\le b_N(\kappa_j)\le c_j.
\]
Taking an occupied tested descendant and a tested ancestor for an arbitrary
occupied intermediate block $I$ of $K$ fine bins proves the finite,
all-scale estimate
\begin{equation}\label{eq:finite-profile-error}
 \left|\log_N|S_i'\cap I|-
 \bigl(s-F_N(s)\bigr)\right|
 \le \max_j(\kappa_{j+1}-\kappa_j)+2\nu+\lambda_N,
 \qquad s=\log_N K.
\end{equation}

For clarity, the passage to arbitrarily many tests is sequential.  At stage
$r$ take a fixed finite mesh of gap at most $1/r$, the first $r$ prescribed
tests and their required joint or domain partitions, and a fixed label
accuracy at most $1/r$.  Let $C_r$ be the reciprocal of the retained mass
fraction guaranteed by the finite construction, including cleaning, and
let $L_r$ be the number of cleaned partitions.  Choose
the resolutions for that stage so large that
$\log(C_rL_r)/\log N\le1/r$ and all other fixed-stage errors are at most
$1/r$.  Let the stage grow sufficiently slowly along the original
resolution sequence.  The resulting refinements are extensive.  Compactness
of the finite profiles $F_N$ gives a uniformly convergent subsequence with
admissible limit $F$.  Equation~\eqref{eq:finite-profile-error} gives
uniformity on all scales.  This proves~(i).

For~(ii), let $s<t$ be fixed.  Uniformity shows that the number of occupied
$s$-blocks in any occupied $t$-block on an original index has exponent
\[
 (t-s)-\bigl(F(t)-F(s)\bigr).
\]
For the restricted and regularized index the corresponding exponent is
$(t-s)-(G(t)-G(s))$.  Its occupied subblocks are a subset of the original
ones.  Comparison proves~\eqref{eq:branching-domination}; rational scales
suffice, followed by continuity.  Let $n'$ be the remaining total index
weight.  Since $n'\le n$, the assumed event-mass retention gives
\[
 n'N^{1-G(1)+o(1)}
 \ge nN^{1-F(1)-e+o(1)}.
\]
This proves the upper bound in~\eqref{eq:total-deficit-loss}; the lower bound
follows from~\eqref{eq:branching-domination} with $s=0,t=1$.  If $e=0$,
all increments of $G-F$ are nonnegative and its total increment is zero,
so $G=F$.

For~(iii), apply the same child-block upper bounds within each derived
fiber.  The sum of its possible starting weights is at most $N^{o(1)}$
times the sum for the original index--starting-block pairs.  The retained
event mass supplies the same lower bound used in the preceding displayed
formula.  Thus the total branching loss is at most the event-loss exponent,
and every subinterval has the inherited upper branching bound.  After
normalizing a logarithmic interval of length $L>0$ to length one, both the
logarithmic loss and the profile deficits are divided by $L$.  No estimate
uniform in a vanishing $L$ is asserted before fixing that interval.
\end{proof}

If earlier tests have already been homogenized, an extensive refinement
can preserve their limiting cell exponents by including their cell and
domain partitions in the cleaning list. The normalized logarithm of every
surviving cell probability then changes by $o(1)$. This is the probability
preservation used for the entropy limits in Section~\ref{sec:entropy};
no such invariance is asserted for tests that have not been homogenized.

\begin{corollary}[Active weight after cleaning]
\label{cor:active-weight-cleaning}
Let $\mathcal A$ be a finite collection of pairs $(i,I)$, where $I$ is a
dyadic temporal block and the blocks assigned to each index $i$ are
disjoint. Let $\mathcal C$ be a partition of $\mathcal A$. For a pair
$(i,I)$, its events are $(i,J)$ with $J\in S_i\cap I$ and have inherited
weight $\omega_i$. Fix $a\in\R$, $\varepsilon\ge0$, and
$0<\vartheta_0,\sigma_N\le1$. Suppose each pair has between
$N^{a-\varepsilon}$ and $N^{a+\varepsilon}$ events.
For $C\in\mathcal C$ put
\[
 M_C=\sum_{(i,I)\in C}\omega_i.
\]
If a refinement retains at least a fraction $\vartheta_0$ of the event
mass assigned to $C$, then its active weight in that cell satisfies
\[
 \sum_{\substack{(i,I)\in C\\S_i'\cap I\ne\varnothing}}\omega_i
 \ge\vartheta_0N^{-2\varepsilon}M_C.
\]
Suppose further that each cell's assigned events are partitioned into at
most $L_N$ pieces, where $L_N\ge1$ is an integer. Discard the pieces of mass less than
$\sigma_N/L_N$ times their original cell's event mass. This discards at
most a fraction $\sigma_N$ of the total event mass. If subsequent cleaning
retains at least a fraction $\vartheta_0$ of each surviving piece's
original event mass, then the active weight on every occupied piece is
at least
\[
 \frac{\vartheta_0\sigma_N}{L_N}N^{-2\varepsilon}M_C.
\]
Here active weight counts each pair $(i,I)$ meeting the piece once.
\end{corollary}

\begin{proof}
The original event mass assigned to $C$ is at least
$N^{a-\varepsilon}M_C$. Each pair contributes at most
$N^{a+\varepsilon}\omega_i$ to its retained event mass, so division by
$N^{a+\varepsilon}$ gives the first bound. The discarded pieces in a
cell have total mass at most $\sigma_N$ times that cell's mass.
An occupied retained piece has mass at least
$\vartheta_0\sigma_N/L_N$ times the original cell's mass; the same
division proves its active-weight bound.
\end{proof}

The cell and piece partitions in this corollary may be included in the
cleaning list. If $L_N=N^{o(1)}$, choose
$\sigma_N\to0$ with $\sigma_N=N^{-o(1)}$; subpower cleaning then
preserves the original active-weight exponent as $\varepsilon\to0$.
For a further extensive selection of events, the cutoff can likewise be
chosen subpower and smaller than its retained fraction. Under a
polynomial restriction the cutoff must lie below that smaller fraction,
and its polynomial loss in the active-weight bound must be retained,
as in~\eqref{eq:total-deficit-loss}.

\begin{convention}[Fixed tests and strict margins]\label{conv:slow-diagonal}
Each argument is first made for a finite list of scales, widths, partitions
and strictly positive tolerances.  Subsequences may be selected so that its
bounded logarithmic parameters converge.  Only after its fixed-parameter
estimates are established are the test lists increased or tolerances sent
to zero, with the resolution chosen large enough for every earlier error.
If a selection is defined by perturbing a parameter, its error tolerance is
sent to zero faster than that fixed perturbation before the latter is
allowed to shrink.  Thus countable tests impose no convergence rate
uniform at all differentiation points or all shrinking scale gaps.
\end{convention}

This convention also explains discarding negligible exceptions.  If a bad
part of a fixed test could have event fraction $N^{-o(1)}$ along a
subsequence, it would itself be an extensive configuration to which the
contradiction applies.  If this is impossible, its fraction is at most
$N^{-c}$ for some fixed $c>0$ eventually.  Finitely many such exceptions can
be removed together; further tests are incorporated only in the subsequent
slow diagonal.  Spatial and angular box tests with continuous parameters
can be covered by finite nets at each fixed resolution, with bounded
thickening.  For all limiting scale exponents it is enough to use increasingly
fine meshes and monotone coverings.

\subsection{Temporal cost under restriction and changes of scale}

Regularization provides the profiles on which the critical inequality is
imposed. We now check how their costs behave under the restrictions,
limits, and changes of scale used to construct equality sequences.

\begin{lemma}[Cost properties]\label{lem:profile-cost}
The cost $\Pi_F$ is additive across logarithmic scale intervals and is
zero on flat intervals. In addition to \eqref{eq:cost-elementary-bounds},
it has the following properties.
\begin{enumerate}[label=(\roman*)]
\item If $F_j\to F$ uniformly, then
$\Pi_F(t)\le\liminf_j\Pi_{F_j}(t)$ for each fixed $t$.
\item If $G$ is obtained as in Lemma~\ref{lem:regularization} with loss
exponent at most $e$, then
\[
 0\le\Pi_G(1)-\Pi_F(1)\le pe.
\]
The same estimate applies on scale subintervals with their corresponding
loss accounting.
\item If a profile on a scale interval $[a,a+L]$ is normalized by
$G(s)=[F(a+Ls)-F(a)]/L$, then
\[
 \Pi_G(1)=\frac1L\int_a^{a+L}P_{p,q}(F'(u))\,du.
\]
Grafting fully occupied time below a cut grafts a flat interval and adds no
cost there.
\end{enumerate}
Moreover $\partial_p\Pi_F(1)=F(1)$.  For $q>0$,
\begin{equation}\label{eq:cost-z-derivative}
 \partial_z\Pi_F(t)=\int_0^t\mathcal D_q(q,F'(u))\,du,
 \qquad
 \mathcal D_q(q,e)=\frac1\eta\int_{1/2-\eta}^{1/2}
                  \one_{\{q<re\}}\,dr.
\end{equation}
\end{lemma}

\begin{proof}
For a partition $0=t_0<\cdots<t_J=t$, Jensen's inequality gives
\[
 \sum_{j=1}^J(t_j-t_{j-1})
 P_{p,q}\!\left(\frac{F(t_j)-F(t_{j-1})}{t_j-t_{j-1}}\right)
 \le\Pi_F(t).
\]
The supremum over refining dyadic partitions equals $\Pi_F(t)$: the
interval averages of the bounded function $F'$ converge to it in $L^1$,
and $P$ is Lipschitz.  Each finite sum is continuous in $F$ for uniform
convergence, proving~(i).  By~\eqref{eq:branching-domination}, $G-F$ is
nondecreasing and Lipschitz, so $G'\ge F'$ almost everywhere.  The slope
bound gives
\[
 0\le\int_0^1[P(G')-P(F')]\le
 p\int_0^1(G'-F')\le pe,
\]
proving~(ii).  A change of variables proves~(iii).

The $p$ derivative is immediate.  For $q>0$ the displayed average for
$\mathcal D_q$ exists also when $e=0$; for $e>0$ its possible equality set
in the $r$ integral has measure zero.  It is continuous in $(q,e)$ with
$q>0$, and bounded between zero and one.  Dominated convergence gives
\eqref{eq:cost-z-derivative}, with the positive sign because $z=-q$.
\end{proof}

\paragraph{Filling occupied blocks.}
Fill every occupied $N^\kappa$-block with all its fine time
bins, keeping the original set of occupied blocks.  In the original
logarithmic units its new profile is
\[
 F^{\mathrm{fill}}(s)=
 \begin{cases}
 0,&0\le s\le\kappa,\\
 F(s)-F(\kappa),&\kappa\le s\le1.
 \end{cases}
\]
Above the cut, count the original occupied cut blocks and multiply by the
full $N^\kappa$ bins in each; below the cut every occupied block is full.
Its cost is precisely $\Pi_F(1)-\Pi_F(\kappa)$ before any further affine
change of the logarithmic scale.  Rebinning must use subdivisions of the
same nested blocks.

\subsection{Uniform bounds for ensembles}

We record precisely how sequence bounds are used on families of local
problems.  Suppose a common upper estimate holds for every bounded-chart
sequence with a specified limiting profile, and every member of an
ensemble has that profile with uniform errors, the same resolution
exponent bounded away from zero, and a common upper bound for its raw plank
parameter.  If the estimate failed uniformly by a fixed power, choosing
one violating member at each resolution would give a forbidden sequence.
Hence the estimate holds uniformly over the ensemble.  Summing its
incidence masses and separately counted support sizes gives the same
upper multiplicity estimate for the aggregate ensemble.  This argument
neither changes the weights nor asserts bounded overlap between supports
in different parents.

The same reasoning is valid for profiles approaching a compact limit when
the upper bound is written with their own costs.  Indeed a fixed-power
failure along profiles $F_j\to F$ gives, by
Lemma~\ref{lem:profile-cost}(i), a failure with cost at least $\Pi_F(1)$
after selecting sufficiently accurate finite realizations.  This would
contradict the all-sequence bound for $F$.  These observations will be used
whenever local charts or representatives are chosen after the original
resolution has been fixed.

Splitting support among distinct charts need not preserve average
multiplicity. In the aggregate just described, support sizes are counted
separately for each chart; a lower multiplicity bound must therefore be
proved for that ensemble in each application.

\subsection{A differentiability point forced by failure}

A failure of the desired bound at $z=0$ will supply an interior parameter
point where both $h$ and its $z$ derivative are positive. The two derivatives
of $h$ at that point will control the temporal deficit and the plank shapes
selected in Section~\ref{sec:stationary}.

\begin{lemma}[Critical parameter point]\label{lem:critical-point}
The function $h(p,z)-z$ is separately nonincreasing in $p$ and $z$.
It is totally differentiable almost everywhere in the interior of its
parameter domain.  If $I$ is an open interval of $p>2$ on which
$h(p,0)>0$, there is a point $(p,z)\in I\times(-1,0)$ at which $h$ is
totally differentiable and
\[
 h(p,z)>0,\qquad h_z(p,z)>0.
\]
At such a point $h_p\le0$ and $h_z\le1$.
\end{lemma}

\begin{proof}
Increasing $p$ increases the cost and therefore cannot increase $h$.
For $z_2=z_1+s$, $s>0$, the dimension parameter increases by $s$.
Starting with a valid candidate $A$ at $z_1$, use $A+s$ at $z_2$.
Its $a$ exponent is unchanged and its $b$ exponent increases by $s$.
Writing the dimension in a subscript temporarily, we have
\[
 \Delta_{A+s,d+s}\ge N^{-s}\Delta_{A,d}.
\]
Also $q$ decreases, so $P_{p,q}$ increases pointwise.  The new right side
of~\eqref{eq:critical-bound} is therefore at least the old one.
The candidate is still in range since $A+s\le d+s-1$.  Thus
\begin{equation}\label{eq:critical-z-monotonicity}
 h(p,z+s)\le h(p,z)+s,
\end{equation}
including admissible endpoints.  This proves the separate monotonicity of
$g(p,z)=h(p,z)-z$.

Here is the passage from this monotonicity to the stated form of
differentiability.  On each compact rectangle in the interior, $g$ is
bounded, and its variations on coordinate lines are bounded by its range.
Integration along those lines shows that both distributional first
derivatives are finite measures.  Thus $g$ is locally of bounded variation.
We use the approximate differentiability theorem for BV functions
\citep[Theorem~3.83]{AFP2000}:
away from a null set there are a value and a linear map giving a first-order
approximation on sets of density one.  Coordinatewise monotonicity upgrades
this approximation to ordinary total differentiability, as follows.

At such a point $x$, fix an error tolerance and a density-one set on which
the approximate linear estimate holds.  For a displacement $u$ of size $r$
and any fixed $\varepsilon>0$, consider small squares of side comparable to
$\varepsilon r$ located coordinatewise just above and just below $x+u$.
They lie within $O(r)$ of $x$.  For sufficiently small $r$ their area is
larger than the bad part of the density-one set in that neighborhood, so
each square contains a comparison point where the linear estimate holds.
Monotonicity bounds $g(x+u)$ between the two comparison values.  The
comparison displacements differ from $u$ by $O(\varepsilon r)$.
The same sandwich first identifies the actual value $g(x)$ with its
approximate value.  Divide the resulting error by $r$, let $r\to0$, and
then let the approximation tolerance and $\varepsilon$ tend to zero.
This is total differentiability at $x$.

For fixed $p\in I$, the decreasing function $z\mapsto g(p,z)$ has only
nonpositive singular variation.  Its ordinary derivative therefore obeys
\[
 g(p,b)-g(p,a)\le\int_a^b g_z(p,z)\,dz.
\]
Adding $b-a$ gives the same inequality for the increment of $h$ and the
integral of $h_z$.  There is no lost upward jump at $z=0$:
\eqref{eq:critical-z-monotonicity} gives
$h(p,-t)\ge h(p,0)-t$.  At the other endpoint,
$0\le h(p,z)\le1+z$ and $h(p,-1)=0$.
Letting $a\downarrow-1$ and $b\uparrow0$ consequently gives
\[
 0<h(p,0)\le\int_{-1}^0h_z(p,z)\,dz.
\]
Fubini and almost-everywhere total differentiability yield an interior
point of total differentiability with $h_z>0$.  At a differentiability point
where $h=0$, nonnegativity of $h$ would force its gradient to vanish.
Hence the selected point has $h>0$ as well.  The derivative signs follow
from the two monotonicity comparisons.
\end{proof}

If $h(p,0)=0$ does not occur for $p>2$ arbitrarily close to two, its
nonnegativity gives an interval immediately above two on which it is
positive.  Choose an open subinterval $I$ with $p_{\min}=\inf I>2$ and fix
\begin{equation}\label{eq:smoothing-choice}
 0<\eta\le\min\left\{\frac1{32},\frac{p_{\min}-2}{16}\right\}.
\end{equation}
This choice is not circular: at $q=0$ the cost and the critical exponent
are independent of $\eta$.  Lemma~\ref{lem:critical-point} then supplies
an interior point for this fixed smoothing.  For the contradiction below
fix that point and put
\begin{equation}\label{eq:critical-point-notation}
 h=h(p,z),\qquad x=d-1-h,
 \qquad y=1+h,\qquad v=-h_p(p,z).
\end{equation}
Thus $h>0$, $h_z>0$, $v\ge0$, $0\le x\le y$, and $x+y=d$.
Strict positivity of $v$ will follow from the positive-deficit equality and
the parameter test in Section~\ref{sec:stationary}.

\subsection{Hybrid equalities and compactness}

We freeze the shape exponent in the plank parameter slightly below $h$
and optimize only the additional power of $N$. The resulting equality
sequences remain sharp under extensive refinements. Keeping the two
exponents distinct also allows the anisotropic argument in
Section~\ref{sec:local} to force a positive temporal deficit before we
pass to an exact critical equality.

\begin{lemma}[Hybrid equality sequences]\label{lem:hybrid-equality}
Fix $p,z,\eta$ with $h=h(p,z)>0$, and let $0\le h_0<h$.
There is a sharp number $H$ with
\[
 h_0<H\le h
\]
such that every sequence with uniform profile $F$ satisfies
\begin{equation}\label{eq:hybrid-upper}
 m\leexp \Delta_{h_0}N^{H+\Pi_F(1)}.
\end{equation}
There is a sequence with a uniform profile attaining equality in
\eqref{eq:hybrid-upper}.

On any extensive refinement of such an equality sequence, followed by
regularization, the profile is unchanged and
\begin{equation}\label{eq:equality-refinement}
 m'\eqexp m,\qquad \Delta'_{h_0}\eqexp\Delta_{h_0}.
\end{equation}
These comparisons refer to refinements of the indexed graph before any
separate geometric change of coordinates.

Finally, suppose $h_{0,j}\uparrow h$ and, for each $j$, a hybrid equality
has profile $F_j$.  After taking a subsequence and sufficiently accurate
realizations, one obtains a sequence with uniform profile $F$ satisfying
\begin{equation}\label{eq:exact-critical-equality}
 m\eqexp\Delta_hN^{h+\Pi_F(1)}.
\end{equation}
If $F_j(1)\ge\alpha_*>0$, then $F(1)\ge\alpha_*$.
\end{lemma}

\begin{proof}
For every uniform-profile sequence form the score
\[
 \limsup_{N\to\infty}
 \left\{\log_N(m/\Delta_{h_0})-\Pi_F(1)\right\},
\]
and let $H$ be the supremum of these scores.  The critical bound at $h$
and $\Delta_h\le\Delta_{h_0}$ give $H\le h$.
If $H\le h_0$, then the bound defining $h$ would hold at $h_0$ for every
sequence, contrary to $h_0<h$.  Thus $H>h_0$.  The supremum definition also
proves~\eqref{eq:hybrid-upper} at $H$ itself.

Choose for each integer $j$ a sequence with profile $F_j$ and score at least
$H-1/j$.  Profiles form a compact set for uniform convergence; pass to a
subsequence with $F_j\to F$.  From the $j$th sequence choose a realization
at a resolution so large that its uniform-profile count error relative to
$F_j$ is at most $1/j$ over all occupied indices and all dyadic time scales,
that all chosen resolutions increase, and that
\[
 \log_N(m/\Delta_{h_0})\ge H+\Pi_{F_j}(1)-2/j.
\]
Any other finite list of prescribed tests is accommodated before this
choice.  The selected diagonal sequence has uniform profile $F$.
Lower semicontinuity of the cost gives
\[
 \liminf\log_N(m/\Delta_{h_0})\ge H+\Pi_F(1).
\]
The universal upper bound~\eqref{eq:hybrid-upper} supplies the reverse
limsup.  Equality follows.  This argument also explains why a downward
jump of the cost at the profile limit cannot cause loss of an extremizer:
it would instead improve the score, which is forbidden by the same sharp
upper bound.  Multiplying all weights by $\Delta_{h_0}^{-1}$ is permitted
if a normalized realization is desired.

For an extensive graph refinement, its incidence mass loses only a
subpower and its support is a subset of the old support.  Hence
$m\leexp m'$.  Its plank masses cannot increase, so
$\Delta'_{h_0}\le\Delta_{h_0}$.  Its profile remains $F$ by
Lemma~\ref{lem:regularization}.  Applying~\eqref{eq:hybrid-upper} to the
refinement sandwiches both quantities:
\[
 m\leexp m'\leexp
 N^{H+\Pi_F(1)}\Delta'_{h_0}
 \le N^{H+\Pi_F(1)}\Delta_{h_0}\eqexp m.
\]
This proves~\eqref{eq:equality-refinement}.

For the final assertion, let $H_j$ denote the corresponding sharp hybrid
power.  By~\eqref{eq:plank-exponent-Lipschitz},
\[
 0\le\log_N(\Delta_{h_{0,j}}/\Delta_h)\le h-h_{0,j},
 \qquad h_{0,j}<H_j\le h.
\]
Thus a sufficiently accurate realization of the $j$th hybrid equality has
\[
 \log_N(m/\Delta_h)-h-\Pi_{F_j}(1)
 =o_j(1),
\]
where the error can be made to tend to zero by first taking that
realization far enough along its sequence.  Take a uniformly convergent
subsequence of $F_j$ and a finite-test diagonal as above.  Lower
semicontinuity gives the lower bound in~\eqref{eq:exact-critical-equality};
the critical bound at $h$ gives its upper bound.  Uniform convergence also
preserves the lower endpoint-deficit bound $F_j(1)\ge\alpha_*$.
\end{proof}

Section~\ref{sec:local} obtains that uniform positive lower bound for the
hybrid deficits from the geometric inputs collected in
Section~\ref{sec:inputs}.  Until then, $H$ and $h_0$ remain distinct in all
local estimates.  Once Lemma~\ref{lem:positive-deficit} is proved,
Lemma~\ref{lem:hybrid-equality} gives an exact equality with positive
deficit at the fixed differentiability point.
